\chapter{Conclusion}
\label{ch:discussion_conclusion}

\section{Summary}
\label{sec:summary}

This work reconstructed turbulent flow fields from sparse point sensors without full-field DNS supervision. PI-CON combines a CfC-DeepONet backbone, isotropic relative-position cross-attention, and an Augmented-Lagrangian-style continuity penalty. On 2D Kolmogorov at $Re=10\,000$ with $K=100$ sensors, it meets the engineering targets across DNS-oracle, DNS-free LES-derived, geometric space-filling, and random placement. LES-derived QR-pivoting is recommended as an informed DNS-free option because it improves over random placement in the evaluated setting, not because it is universally optimal: space-filling and the DNS oracle give lower KE error, whereas LES-derived placement gives lower pointwise error (Table~\ref{tab:placement_strategy_new}).

The three engineering success criteria of \S\ref{sec:research_questions} are met:
\begin{itemize}
    \item \textbf{(O1) An accurate and fast reconstructor.} The energy-dominant low-wavenumber band reaches single-digit relative error within the $10\%$ target (Fig.~\ref{fig:band_err}), with a $2\times2$ ablation confirming the architectural components under the final PI-CON protocol; once trained, the operator answers any space--time query in a single forward pass and reconstructs a full trajectory $20\times$ faster than re-solving the reference DNS on the hardware reported, meeting the $\ge 5\times$ criterion (Table~\ref{tab:inference_cost}), though a single full-field query is not real-time on CPU or MPS in the current implementation.
    \item \textbf{(O2) Sensor count sets the achievable resolution.} A divergence-free SVD and a $K$-scaling sweep ($K\in\{100,200,400\}$) attribute residual mid/high-band error to sensor information capacity, with the kinetic-energy- and vorticity-derived cutoffs bracketing the sensor-count sampling band edge, identifying sensor-density scaling as the productive direction (\S\ref{sec:info_limit}).
    \item \textbf{(O3) Placement and noise set reliability.} Four placement strategies (DNS-oracle, DNS-free LES, geometric space-filling, and random) all keep KE error within the engineering target. LES-derived QR-pivoting improves over random placement and is therefore recommended as an informed DNS-free option, while the metric-dependent ranking precludes a universal best-placement claim. Placement contributes a variance source $6.4\times$ training stochasticity; additive sensor noise up to $10\%$ keeps the reconstruction within the target in the noise-stress test.
\end{itemize}

\paragraph{Contributions.}
\begin{enumerate}
    \item \textbf{PI-CON architecture for engineering sparse-sensor reconstruction.} Under sensor-only-with-physics supervision, PI-CON attains KE MAPE $5.71 \pm 0.12\%$ ($n=5$ seeds) from $K=100$ sensors at $Re=10\,000$ (Table~\ref{tab:main_metrics}) and reconstructs any query point without re-solving the governing equations. The matched $2\times2$ ablation identifies cross-attention as the dominant architectural lever, with CfC useful through its interaction with the query-conditional readout.
    \item \textbf{A systematic sensing-configuration study.} Sensor count, placement, and noise are separated experimentally. A divergence-free SVD and $K\in\{100,200,400\}$ sweep show that count sets resolution (Table~\ref{tab:k_scaling_nyquist}); placement sets reliability, with placement variance $6.4\times$ training stochasticity and no single strategy dominating all metrics (Table~\ref{tab:placement_strategy_new}); and additive noise up to $10\%$ leaves reconstruction within the target. Sensor-density scaling, not architecture, is the productive direction for higher fidelity.
    \item \textbf{Cross-Reynolds feasibility under a single architecture.} At $Re=10^6$ with $K=200$ sensors the same architecture--AL recipe attains KE MAPE $6.10\%$ (Table~\ref{tab:re1e6}), showing feasibility over two orders of magnitude in Reynolds number. The result is single-seed and should be treated as an extension, not a multi-seed benchmark.
\end{enumerate}

\section{Engineering Implications}
\label{sec:engineering_impl}

Five implications follow:

\begin{itemize}
    \item \textbf{Energy-dominant band suffices for engineering use.} The low-wavenumber band is reconstructed at single-digit relative error from a DNS-free LES-pivot pipeline, sufficient for KE-based control, mean-flow monitoring, and phase-locked actuation.
    \item \textbf{Mid/high-band saturation has a quantified bound.} The near-$100\%$ higher-band error tracks the sensor-information limit, and $K$-scaling shows the effective bandwidth expanding with the predicted $\sqrt{K}$ scaling (\S\ref{sec:info_limit}).
    \item \textbf{Active incompressibility constraint across Reynolds numbers.} The Augmented-Lagrangian continuity constraint drives the reconstruction's divergence to the finite-difference floor of its resolved bandwidth at both reported Reynolds numbers (\S\ref{sec:sub_dns_div}), with no configuration switching required. Incompressibility-sensitive downstream tasks should validate against their own tolerance, since the reconstruction is band-limited near the sensor-count sampling band edge.
    \item \textbf{Streaming-friendly deployment.} Filtering mode reads no future sensors and the CfC formulation decouples query rate from sensor sampling; training uses no DNS full field, so deployment requires only the sensor stream.
    \item \textbf{Bounded sensitivity to observation noise.} Across the tested additive-noise range, KE remains below the $10\%$ engineering threshold, supporting robustness to modest observation noise while leaving realistic sensor-error models open.
\end{itemize}

\section{Limitations}
\label{sec:limitations}

\begin{itemize}
    \item \textbf{Observation noise.} The $1$--$10\%$ Gaussian noise sweep uses the final PI-CON protocol with $n=5$ seeds per noise level, but it still covers only additive per-channel Gaussian perturbations and excludes sensor-specific bias, drift, dropout, correlated channel noise, and calibration errors.
    \item \textbf{Geometry scope.} The main study uses a simple, obstacle-free, doubly periodic two-dimensional domain. A single non-periodic case (a cylinder wake with an explicit inflow boundary-condition loss) is validated as a preliminary extension in Appendix~\ref{app:cylinder}, but a systematic treatment of wall-bounded and obstacle-laden geometries (airfoil, internal flow) remains open.
    \item \textbf{Dimensionality.} All reported results are two-dimensional. Two-dimensional turbulence carries an inverse energy cascade and no vortex stretching, so the spectral content the reconstruction must recover differs in kind from three-dimensional turbulence. The sensor-information analysis of \S\ref{sec:info_limit} is likewise derived for a two-dimensional divergence-free basis, and its scaling does not transfer to three dimensions without rederivation.
    \item \textbf{Pressure field.} Only the velocity components $(u,v)$ are reconstructed under supervision. Pressure appears as a model-internal variable coupled to the loss solely through $\nabla p$ in the momentum residual, leaving both an additive gauge and the pressure magnitude underdetermined. The resulting internal pressure departs from the DNS reference in both gradient and gauge-removed value, and the departure is independent of the AL recipe (Appendix~\ref{app:pressure_limit}). Pressure-derived quantities --- surface loading, cavitation margin, and aeroacoustic source terms --- are therefore outside the validated scope.
    \item \textbf{Forcing form.} A single sustained body force $f_x = A\sin(2\pi k_f y)$ is used throughout. The architecture is forcing-agnostic, but wall-driven channel flow and inflow-driven wakes/jets require case-by-case revalidation.
    \item \textbf{Temporal sampling.} The CfC branch is adopted because a deployable temporal model must tolerate the irregular, unsynchronised timestamps of field sensors (\S\ref{subsec:cfc_encoder}), yet the controlled Kolmogorov benchmark samples uniformly at $0.025$~s. The irregular-clock capability is therefore carried by the formulation rather than exercised by the evidence reported here, and the $2\times2$ ablation (\S\ref{sec:baseline_comparison}) measures the CfC branch only on a fixed clock, where the time gap entering the cell is constant; this is the regime in which the branch is least able to pay for its aggregation cost. Whether the capability holds under non-uniform timestamps remains an open follow-up (\S\ref{sec:future_work}).
    \item \textbf{Acceptance metrics.} Primary KE, $u/v$ relative $L_2$, divergence, mean-profile, Reynolds-shear-stress, and enstrophy/dissipation diagnostics are covered; higher-order Reynolds-stress moments and full turbulent-kinetic-energy budget closure are not.
    \item \textbf{Per-case fitting and single-seed cross-$Re$.} PI-CON is fit once per case by design because high-Reynolds chaos ties each realisation to its conditions. Cross-case generality is not demonstrated, and the $Re=10^6$ extension is single-seed with unquantified seed variance.
    \item \textbf{CFD-rigour gaps.} The DNS passes Pope-2000 ($k_{\max}\eta=1.91$) but $T/t_{\rm eddy}=2.51$ gives a short statistical window.
    \item \textbf{Diagnostic boundaries.} Forcing parameters ($A$, $k_f$) were not recovered at $K=100$ sensors under the tested initialisation despite accurate reconstruction (\S\ref{sec:forcing_identifiability}); structural identifiability is not established. Multi-constraint AL remains a recipe diagnostic rather than a main conclusion: the final-protocol rerun gives mixed single-seed results, so the thesis keeps the continuity-only AL recipe, supported by the $n=5$ main baseline (\S\ref{sec:multi_al_negative}).
\end{itemize}

\section{Future Work}
\label{sec:future_work}

Nine directions follow from the limitations:
\begin{enumerate}
    \item \textbf{Sensor-budget scaling beyond the current curve.} Extend the $K=100/200/400$ sweep with $K=50$ and a matched training budget at $K=400$ to stress-test the agreement between the sensor-count sampling band edge and the effective reconstruction bandwidth, and to separate sensor bandwidth from compute-budget compromises.
    \item \textbf{Cross-Reynolds multi-seed confirmation.} The $Re=10^6$ extension (Table~\ref{tab:re1e6}) is single-seed; an $n=3$ multi-seed run quantifies its seed variance against the $Re=10^4$ baseline. A higher-Reynolds test ($Re \ge 10^7$) and sensor-density scaling at $Re=10^6$ ($K=400$) extend the cross-$Re$ envelope.
    \item \textbf{Realistic sensor-error model.} Go beyond additive Gaussian noise to bias, drift, dropout, correlated channel noise, and calibration uncertainty; assess whether sensor denoising or mean-enforced losses such as that of Mo and Magri~\cite{Mo2025PRF} are needed.
    \item \textbf{Irregular-clock validation.} Re-run the reconstruction on non-uniform sensor timestamps --- dropped samples and per-group cadences --- against a discrete recurrent baseline that receives the time gap as an input feature. Masking the time axis of the existing $201$-frame trajectories supplies such a stream without a new reference solve, making this the cheapest of the directions listed here, and it tests the continuous-time property on which the branch is selected (\S\ref{subsec:cfc_encoder}). A CfC main effect that stays positive under an irregular clock, or a CfC$\times$cross-attention interaction that does not widen, would place the continuous-time argument outside the present problem class rather than merely outside the present benchmark.
    \item \textbf{Wall-bounded and obstacle geometries with multi-modal fusion.} Extend the preliminary cylinder-wake validation (Appendix~\ref{app:cylinder}) to a systematic treatment of non-periodic, wall-bounded, and obstacle-laden two-dimensional flows (airfoil, backward-facing step) under explicit boundary-condition losses, incorporating pressure taps or temperature sensors through an extended per-channel embedding $\mathbf{e}_c$.
    \item \textbf{Pressure-field reconstruction.} The pressure gap is structural rather than a tuning failure: the momentum residual constrains only $\nabla p$, so the loss is invariant to an additive gauge $p\to p+C(t)$, and any smooth pattern partially cancelling $(\mathbf{u}\cdot\nabla)\mathbf{u}$ is admissible (Appendix~\ref{app:pressure_limit}). Three routes address it, in increasing order of architectural commitment.

    The first adds a sparse pressure-gauge channel. Point pressure taps are among the cheapest and most standard instruments in an engineering test rig --- far easier to install than interior velocity probes --- so a handful of taps is consistent with the deployable-supervision constraint of \S\ref{sec:engineering_constraint}. Entering them through the per-channel embedding $\mathbf{e}_c$ requires no architectural change: $K_p$ tap readings fix the gauge and anchor the low-order pressure modes, and the same count/placement/noise study applied here to velocity (\S\ref{sec:info_limit}) then determines how many taps are needed and where they belong. The wall-tap variant, where taps sit only on a boundary rather than in the interior, is the configuration most representative of practice and the one most likely to reveal an observability limit specific to pressure.

    The second reparametrises pressure out of the output. Taking the divergence of the momentum equation gives the pressure-Poisson equation $\nabla^2 p = -\nabla\cdot\big[(\mathbf{u}\cdot\nabla)\mathbf{u}\big]$, whose solution makes $p$ a deterministic functional of the reconstructed velocity rather than a free network output, removing the underdetermination by construction. The cost is a transfer of burden: the right-hand side depends on velocity gradients, and the reconstruction is band-limited near the sensor-count sampling band edge, so the sharpest test of this route is whether the recovered pressure degrades faster than the velocity field it is derived from.

    The third keeps $p$ as an output but adds the pressure-Poisson residual as a second Augmented-Lagrangian constraint alongside continuity, a softer intervention that reuses the existing multiplier machinery (\S\ref{sec:multi_al_negative}) and can be compared against the mixed evidence already reported for multi-constraint AL. Whichever route is taken, the evaluation should report the pressure-gradient and gauge-removed errors of Appendix~\ref{app:pressure_limit} against the same DNS reference, so that progress is measured on the metric on which the present scope limit was established.
    \item \textbf{Three-dimensional wall-bounded turbulence.} Turbulent channel flow at a friction Reynolds number is the natural target beyond the present scope, and resolving it is a separate research programme rather than an extension of the experiments reported here. Three obstacles are specific to it. Wall-normal inhomogeneity removes the periodicity that the isotropic relative-position cross-attention and the periodic collocation sampling both assume, so placement must become wall-normal-aware and near-wall gradients must be represented under a stretched coordinate. The sensor-information argument also changes character: the two-dimensional sampling band edge $\sqrt{K/\pi}$ derives from $K$ sensors spread over an area, whereas an equivalent bandwidth in a volume scales as $K^{1/3}$, so matching the present resolved bandwidth would demand a sensor budget beyond what interior probes can supply --- pushing the problem towards wall-only measurement (wall shear stress and wall pressure) combined with inner/outer scaling priors, and coupling this direction to the pressure-channel work above. Vortex stretching and the forward energy cascade further alter the spectral content that physics residuals must constrain, so the Reynolds-number envelope established in two dimensions carries no guarantee. Reference-data and training cost scale accordingly, which makes this the direction where the training--simulation crossover below is decided.
    \item \textbf{4D-Var assimilation.} Condition a forward solver on the continuous sensor stream via 4D-Var and compare the assimilated trajectory against the operator reconstruction on matched sensor configurations, separating what the physics prior contributes from what the learned operator contributes.
    \item \textbf{Cross-case generalisation and the training--simulation crossover.} Train a single model on multi-Reynolds, multi-trajectory data and test whether it generalises to unseen cases, quantifying the accuracy cost of amortising per-case fitting. In regimes where one high-fidelity solve is prohibitive (3D, multi-physics, or substantially higher Reynolds number), measure whether the fixed offline training cost falls below a single forward simulation, the condition under which operator learning would undercut direct simulation end-to-end.
\end{enumerate}
